\begin{resumo}
Este projeto tem como objetivo desenvolver uma interface web para monitoramento e controle de um robô autônomo navegando em labirintos, oferecendo ao operador a visualização em tempo real das leituras dos sensores Time of Flight (ToF) frontais e laterais, com auditoria da precisão na detecção de paredes. Entre os principais requisitos funcionais, destacam-se a autenticação segura de operadores por meio de tela de login com validação de credenciais e controle de sessão; a exibição dinâmica das paredes detectadas na matriz visual assim que os valores brutos dos sensores estabilizarem abaixo do limiar configurado; e a comunicação eficiente entre o frontend e o firmware do robô para recebimento contínuo dos dados de distância. O projeto foi desenvolvido no contexto do Projeto Integrador de Engenharia 1 da Faculdade UnB Gama – FCTE, consolidando conhecimentos de programação web, usabilidade e sistemas embarcados em uma aplicação prática de robótica móvel.

\vspace{\onelineskip}
\noindent
\textbf{Palavras-chave}: Projeto Integrador de Engenharia 1. Faculdade UnB Gama -- FCTE. Robótica Móvel. Interface Web. Sensores ToF.
\end{resumo}

% \begin{resumo}
%  O resumo deve ressaltar o objetivo, o método, os resultados e as conclusões 
%  do documento. A ordem e a extensão
%  destes itens dependem do tipo de resumo (informativo ou indicativo) e do
%  tratamento que cada item recebe no documento original. O resumo deve ser
%  precedido da referência do documento, com exceção do resumo inserido no
%  próprio documento. (\ldots) As palavras-chave devem figurar logo abaixo do
%  resumo, antecedidas da expressão Palavras-chave:, separadas entre si por
%  ponto e finalizadas também por ponto. O texto pode conter no mínimo 150 e 
%  no máximo 500 palavras, é aconselhável que sejam utilizadas 200 palavras. 
%  E não se separa o texto do resumo em parágrafos.

%  \vspace{\onelineskip}
    
%  \noindent
%  \textbf{Palavras-chave}: latex. abntex. editoração de texto.
% \end{resumo}
